\documentclass{beamer}
\usetheme{Madrid}
\usecolortheme{default}
\usepackage[utf8]{inputenc}
\usepackage{amsmath}
\usepackage{amsfonts}
\usepackage{amssymb}
\usepackage{tikz}

\title{CRISC Domain A - 100 More Impossible Hard Questions}
\subtitle{Governance - Expert Level Part 2}
\author{CRISC Mastery Series}
\date{\today}

\begin{document}
	
	\begin{frame}
		\titlepage
	\end{frame}
	
	\begin{frame}{Warning}
		\begin{alertblock}{Extreme Difficulty}
			These questions test the deepest governance judgment.\\
			Every option will look plausible.\\
			You must think like a board advisor, not a technician.
		\end{alertblock}
	\end{frame}
	
	\begin{frame}{Outline}
		\tableofcontents
	\end{frame}
	
	% ============================================================
	\section{1. Strategic Alignment and Governance}
	% ============================================================
	
	\begin{frame}{Q101: Strategy Cascade}
		\begin{block}{Question}
			The CEO announces a digital transformation strategy. The risk practitioner must translate this into IT risk priorities. What is the BEST FIRST step?
		\end{block}
		\begin{enumerate}[(A)]
			\item Purchase new tools
			\item Identify business objectives, drivers, and dependencies introduced by the strategy
			\item Hire more analysts
			\item Approve more projects
		\end{enumerate}
	\end{frame}
	
	\begin{frame}{Q101: Explanation}
		\begin{block}{Correct Answer: B}
			\textbf{Strategy defines the risk agenda:}
			\begin{itemize}
				\item Strategy reveals new dependencies and objectives
				\item Objectives anchor risk priorities
				\item Tools and headcount follow the priorities
			\end{itemize}
		\end{block}
		\begin{alertblock}{Critical Trap}
			Strategy defines the risk agenda, not the reverse.
		\end{alertblock}
	\end{frame}
	
	\begin{frame}{Q102: Board Knowledge of IT}
		\begin{block}{Question}
			The board admits it does not understand a major IT risk presented to it. What governance action is BEST?
		\end{block}
		\begin{enumerate}[(A)]
			\item Approve and move on
			\item Request a focused briefing and independent verification before decision
			\item Postpone indefinitely
			\item Delegate to the IT vendor
		\end{enumerate}
	\end{frame}
	
	\begin{frame}{Q102: Explanation}
		\begin{block}{Correct Answer: B}
			\textbf{Understanding precedes decision:}
			\begin{itemize}
				\item Understanding precedes decision
				\item Independent verification reduces blind spots
				\item Postponement delays important action
			\end{itemize}
		\end{block}
		\begin{alertblock}{Critical Trap}
			Boards must understand before deciding.
		\end{alertblock}
	\end{frame}
	
	\begin{frame}{Q103: Ethical Judgment Escalation}
		\begin{block}{Question}
			An executive directs a team to hide a security incident from the board. What should the practitioner do?
		\end{block}
		\begin{enumerate}[(A)]
			\item Comply silently
			\item Escalate through ethics, audit, or legal channels
			\item Ignore
			\item Delay reporting
		\end{enumerate}
	\end{frame}
	
	\begin{frame}{Q103: Explanation}
		\begin{block}{Correct Answer: B}
			\textbf{Concealment violates governance:}
			\begin{itemize}
				\item Concealment violates governance
				\item Channels exist for escalation
				\item Silence compounds exposure
			\end{itemize}
		\end{block}
		\begin{alertblock}{Critical Trap}
			Concealment violates the Code of Ethics.
		\end{alertblock}
	\end{frame}
	
	\begin{frame}{Q104: Governance of Decisions}
		\begin{block}{Question}
			An IT decision changes the enterprise risk profile materially. What governance principle MOST applies?
		\end{block}
		\begin{enumerate}[(A)]
			\item IT can decide independently
			\item Material changes require escalation and approval per governance thresholds
			\item Vendors decide
			\item The change can proceed
		\end{enumerate}
	\end{frame}
	
	\begin{frame}{Q104: Explanation}
		\begin{block}{Correct Answer: B}
			\textbf{Material change needs governance approval:}
			\begin{itemize}
				\item Material changes need authority
				\item Thresholds define escalation
				\item Silent change loses oversight
			\end{itemize}
		\end{block}
		\begin{alertblock}{Critical Trap}
			Material change needs governance approval.
		\end{alertblock}
	\end{frame}
	
	\begin{frame}{Q105: Policy Cascade}
		\begin{block}{Question}
			The board approves a new high-level policy. What is the NEXT governance step?
		\end{block}
		\begin{enumerate}[(A)]
			\item Archive the policy
			\item Cascade into standards and procedures with owners and timelines
			\item Notify regulators
			\item Approve budget
		\end{enumerate}
	\end{frame}
	
	\begin{frame}{Q105: Explanation}
		\begin{block}{Correct Answer: B}
			\textbf{Policy without cascade is a document:}
			\begin{itemize}
				\item Policy requires implementation
				\item Standards and procedures operationalize intent
				\item Owners and timelines drive action
			\end{itemize}
		\end{block}
		\begin{alertblock}{Critical Trap}
			Policy without cascade is a document, not a control.
		\end{alertblock}
	\end{frame}
	
	\begin{frame}{Q106: Risk Appetite Revisited}
		\begin{block}{Question}
			The business strategy shifts from stability to aggressive growth. What should happen to risk appetite?
		\end{block}
		\begin{enumerate}[(A)]
			\item Remain unchanged
			\item Be reconsidered and re-approved to reflect the new strategy
			\item Decrease
			\item Increase without review
		\end{enumerate}
	\end{frame}
	
	\begin{frame}{Q106: Explanation}
		\begin{block}{Correct Answer: B}
			\textbf{Appetite follows strategy:}
			\begin{itemize}
				\item Appetite follows strategy
				\item Board re-approves appetite
				\item Changes must be deliberate
			\end{itemize}
		\end{block}
		\begin{alertblock}{Critical Trap}
			Appetite follows strategy, with board approval.
		\end{alertblock}
	\end{frame}
	
	\begin{frame}{Q107: Management Versus Board}
		\begin{block}{Question}
			Who decides whether to accept a risk that exceeds the stated appetite?
		\end{block}
		\begin{enumerate}[(A)]
			\item The risk practitioner
			\item The board of directors or its delegated authority per policy
			\item The IT manager
			\item The vendor
		\end{enumerate}
	\end{frame}
	
	\begin{frame}{Q107: Explanation}
		\begin{block}{Correct Answer: B}
			\textbf{Above-appetite acceptance is board-owned:}
			\begin{itemize}
				\item Accepting above appetite is a board-level decision
				\item Delegation must be defined
				\item Practitioners advise, not decide
			\end{itemize}
		\end{block}
		\begin{alertblock}{Critical Trap}
			Above-appetite acceptance is board-owned.
		\end{alertblock}
	\end{frame}
	
	\begin{frame}{Q108: Board Committee Structure}
		\begin{block}{Question}
			A board creates a technology risk committee. What is the PRIMARY benefit?
		\end{block}
		\begin{enumerate}[(A)]
			\item More meetings
			\item Focused expertise and oversight on technology risk
			\item Cheaper operations
			\item Avoidance of audits
		\end{enumerate}
	\end{frame}
	
	\begin{frame}{Q108: Explanation}
		\begin{block}{Correct Answer: B}
			\textbf{Committees deepen expertise:}
			\begin{itemize}
				\item Committees deepen expertise
				\item Focus improves decisions
				\item Costs and audits are not the goal
			\end{itemize}
		\end{block}
		\begin{alertblock}{Critical Trap}
			Committees deepen expertise and focus.
		\end{alertblock}
	\end{frame}
	
	\begin{frame}{Q109: Risk Committee Authority}
		\begin{block}{Question}
			A risk committee debates risks but has no authority to require action. What is the MAIN gap?
		\end{block}
		\begin{enumerate}[(A)]
			\item Too many members
			\item Missing authority and escalation paths to enforce actions
			\item Too much data
			\item Too few vendors
		\end{enumerate}
	\end{frame}
	
	\begin{frame}{Q109: Explanation}
		\begin{block}{Correct Answer: B}
			\textbf{Debate without authority stalls action:}
			\begin{itemize}
				\item Committees need authority or escalation
				\item Debate without power delays action
				\item Defined authority supports effectiveness
			\end{itemize}
		\end{block}
		\begin{alertblock}{Critical Trap}
			Debate without authority stalls action.
		\end{alertblock}
	\end{frame}
	
	\begin{frame}{Q110: Role Clarity}
		\begin{block}{Question}
			The CISO and the CRO both claim ownership of IT risk reporting. What is the BEST resolution?
		\end{block}
		\begin{enumerate}[(A)]
			\item The CISO reports
			\item Define responsibilities through a RACI or governance policy; the CRO typically owns enterprise risk reporting
			\item The CIO reports
			\item Vendor reports
		\end{enumerate}
	\end{frame}
	
	\begin{frame}{Q110: Explanation}
		\begin{block}{Correct Answer: B}
			\textbf{Overlaps are resolved with clarity:}
			\begin{itemize}
				\item Overlap creates conflict
				\item RACI clarifies roles
				\item Enterprise risk reporting usually sits with the CRO
			\end{itemize}
		\end{block}
		\begin{alertblock}{Critical Trap}
			Overlaps are resolved with clarity, not hierarchy.
		\end{alertblock}
	\end{frame}
	
	\begin{frame}{Q111: M and A Governance}
		\begin{block}{Question}
			An acquisition introduces unknown IT risk. What governance decision is MOST urgent?
		\end{block}
		\begin{enumerate}[(A)]
			\item Approve branding changes
			\item Assign ownership and authorize an assessment before integration
			\item Increase marketing budget
			\item Cancel the acquisition
		\end{enumerate}
	\end{frame}
	
	\begin{frame}{Q111: Explanation}
		\begin{block}{Correct Answer: B}
			\textbf{Acquisitions need owners and assessment:}
			\begin{itemize}
				\item Ownership enables action
				\item Assessment prevents surprises
				\item Cancellation needs evidence
			\end{itemize}
		\end{block}
		\begin{alertblock}{Critical Trap}
			Acquisitions need owners and assessment first.
		\end{alertblock}
	\end{frame}
	
	\begin{frame}{Q112: Board Oversight of Outages}
		\begin{block}{Question}
			Repeated outages are reported to the board without corrective action. What is the BEST board response?
		\end{block}
		\begin{enumerate}[(A)]
			\item Approve the status quo
			\item Require a remediation plan with owners and timelines
			\item Ignore the outages
			\item Reduce monitoring
		\end{enumerate}
	\end{frame}
	
	\begin{frame}{Q112: Explanation}
		\begin{block}{Correct Answer: B}
			\textbf{Repeated failures demand remediation:}
			\begin{itemize}
				\item Boards require action
				\item Owners and timelines create accountability
				\item Ignoring normalizes failure
			\end{itemize}
		\end{block}
		\begin{alertblock}{Critical Trap}
			Repeated failures demand remediation plans.
		\end{alertblock}
	\end{frame}
	
	\begin{frame}{Q113: Executive Removal of Controls}
		\begin{block}{Question}
			An executive orders removal of a control to speed a project. What governance principle MOST applies?
		\end{block}
		\begin{enumerate}[(A)]
			\item Executive authority overrides all controls
			\item Changes to controls require risk assessment and approval per policy
			\item The project should proceed
			\item Vendors should decide
		\end{enumerate}
	\end{frame}
	
	\begin{frame}{Q113: Explanation}
		\begin{block}{Correct Answer: B}
			\textbf{Removing controls needs assessment:}
			\begin{itemize}
				\item Controls exist for reasons
				\item Removal requires assessment and authority
				\item Speed is not a reason to bypass governance
			\end{itemize}
		\end{block}
		\begin{alertblock}{Critical Trap}
			Removing controls needs assessment and authority.
		\end{alertblock}
	\end{frame}
	
	\begin{frame}{Q114: Decision Rights Matrix}
		\begin{block}{Question}
			The organization lacks defined decision rights for IT risk. What is the MOST likely consequence?
		\end{block}
		\begin{enumerate}[(A)]
			\item Faster decisions
			\item Confusion, delays, and unassigned accountability
			\item Lower cost
			\item Better audits
		\end{enumerate}
	\end{frame}
	
	\begin{frame}{Q114: Explanation}
		\begin{block}{Correct Answer: B}
			\textbf{Undefined rights create drift:}
			\begin{itemize}
				\item Without rights, decisions drift or stall
				\item Roles become ambiguous
				\item Governance requires clarity
			\end{itemize}
		\end{block}
		\begin{alertblock}{Critical Trap}
			Undefined rights create drift and confusion.
		\end{alertblock}
	\end{frame}
	
	\begin{frame}{Q115: Board Accountability}
		\begin{block}{Question}
			A major failure occurs after the board ignored warnings from management. Who is accountable?
		\end{block}
		\begin{enumerate}[(A)]
			\item Management only
			\item The board and management, per their defined roles
			\item The auditor
			\item The vendor
		\end{enumerate}
	\end{frame}
	
	\begin{frame}{Q115: Explanation}
		\begin{block}{Correct Answer: B}
			\textbf{Accountability follows oversight roles:}
			\begin{itemize}
				\item Boards oversee
				\item Management operates
				\item Both hold defined accountabilities
			\end{itemize}
		\end{block}
		\begin{alertblock}{Critical Trap}
			Accountability follows defined oversight roles.
		\end{alertblock}
	\end{frame}
	
	% ============================================================
	\section{2. Legal, Regulatory and Ethical Requirements}
	% ============================================================
	
	\begin{frame}{Q116: Data Breach Notification}
		\begin{block}{Question}
			A breach affects customers in multiple countries with different notification rules. What should governance require?
		\end{block}
		\begin{enumerate}[(A)]
			\item Notify only one country
			\item Use a consistent notification process that satisfies all applicable requirements
			\item Delay until resolution
			\item Avoid notifications
		\end{enumerate}
	\end{frame}
	
	\begin{frame}{Q116: Explanation}
		\begin{block}{Correct Answer: B}
			\textbf{Notification must satisfy every jurisdiction:}
			\begin{itemize}
				\item Consistency supports compliance
				\item Multiple rules require careful design
				\item Delay violates laws
			\end{itemize}
		\end{block}
		\begin{alertblock}{Critical Trap}
			Breach notification must satisfy every jurisdiction.
		\end{alertblock}
	\end{frame}
	
	\begin{frame}{Q117: Data Sovereignty}
		\begin{block}{Question}
			A regulation requires data to remain in-country. The vendor stores data abroad. What governance action is BEST?
		\end{block}
		\begin{enumerate}[(A)]
			\item Accept the vendor's statement
			\item Reassess the vendor relationship and design controls to meet sovereignty requirements
			\item Ignore the regulation
			\item Expand the vendor's role
		\end{enumerate}
	\end{frame}
	
	\begin{frame}{Q117: Explanation}
		\begin{block}{Correct Answer: B}
			\textbf{Sovereignty overrides convenience:}
			\begin{itemize}
				\item Sovereignty is a hard constraint
				\item Controls must meet the requirement
				\item Ignoring the law is not an option
			\end{itemize}
		\end{block}
		\begin{alertblock}{Critical Trap}
			Sovereignty overrides vendor convenience.
		\end{alertblock}
	\end{frame}
	
	\begin{frame}{Q118: Ethical Sourcing}
		\begin{block}{Question}
			A supplier violates human rights standards. What governance response is BEST?
		\end{block}
		\begin{enumerate}[(A)]
			\item Ignore it
			\item Assess exposure, require remediation, or exit the relationship per policy
			\item Increase orders
			\item Hide it
		\end{enumerate}
	\end{frame}
	
	\begin{frame}{Q118: Explanation}
		\begin{block}{Correct Answer: B}
			\textbf{Ethical sourcing is a governance obligation:}
			\begin{itemize}
				\item Ethics are part of governance
				\item Options include remediation or exit
				\item Ignoring damages reputation and values
			\end{itemize}
		\end{block}
		\begin{alertblock}{Critical Trap}
			Ethical sourcing is a governance obligation.
		\end{alertblock}
	\end{frame}
	
	\begin{frame}{Q119: Conflicts with Law}
		\begin{block}{Question}
			An internal policy conflicts with local law. What is the BEST governance action?
		\end{block}
		\begin{enumerate}[(A)]
			\item Follow policy
			\item Follow the law and escalate to update the policy
			\item Ignore both
			\item Delegate to the vendor
		\end{enumerate}
	\end{frame}
	
	\begin{frame}{Q119: Explanation}
		\begin{block}{Correct Answer: B}
			\textbf{Law outranks internal policy:}
			\begin{itemize}
				\item Law takes precedence over internal policy
				\item Policy must be updated
				\item Silence endangers the enterprise
			\end{itemize}
		\end{block}
		\begin{alertblock}{Critical Trap}
			Law outranks internal policy.
		\end{alertblock}
	\end{frame}
	
	\begin{frame}{Q120: Professional Ethics}
		\begin{block}{Question}
			The practitioner is asked to sign a report they believe is materially misleading. What should they do?
		\end{block}
		\begin{enumerate}[(A)]
			\item Sign it
			\item Refuse, document concerns, and escalate per the Code of Ethics
			\item Modify the report silently
			\item Resign without action
		\end{enumerate}
	\end{frame}
	
	\begin{frame}{Q120: Explanation}
		\begin{block}{Correct Answer: B}
			\textbf{The Code of Ethics requires honesty:}
			\begin{itemize}
				\item Honesty is a professional duty
				\item Documentation preserves evidence
				\item Escalation is the proper path
			\end{itemize}
		\end{block}
		\begin{alertblock}{Critical Trap}
			The Code of Ethics requires honesty.
		\end{alertblock}
	\end{frame}
	
	\begin{frame}{Q121: Confidentiality Obligations}
		\begin{block}{Question}
			A client's confidential data is requested by a third party. What governs disclosure?
		\end{block}
		\begin{enumerate}[(A)]
			\item Curiosity
			\item Legal obligation or authorized request under contract and law
			\item Convenience
			\item No rules apply
		\end{enumerate}
	\end{frame}
	
	\begin{frame}{Q121: Explanation}
		\begin{block}{Correct Answer: B}
			\textbf{Disclosure requires legal authority:}
			\begin{itemize}
				\item Disclosure needs legal authority
				\item Contracts define permitted sharing
				\item Casual disclosure breaches confidentiality
			\end{itemize}
		\end{block}
		\begin{alertblock}{Critical Trap}
			Disclosure requires legal authority.
		\end{alertblock}
	\end{frame}
	
	\begin{frame}{Q122: Cross-Border Transfer}
		\begin{block}{Question}
			Data must be transferred between jurisdictions. What governance step is MOST important?
		\end{block}
		\begin{enumerate}[(A)]
			\item Transfer without checks
			\item Confirm lawful transfer mechanisms and controls are in place
			\item Delay indefinitely
			\item Use the cheapest method
		\end{enumerate}
	\end{frame}
	
	\begin{frame}{Q122: Explanation}
		\begin{block}{Correct Answer: B}
			\textbf{Lawful mechanisms precede transfers:}
			\begin{itemize}
				\item Lawful mechanisms are required
				\item Controls ensure protection
				\item Cost is not the primary factor
			\end{itemize}
		\end{block}
		\begin{alertblock}{Critical Trap}
			Lawful mechanisms precede transfers.
		\end{alertblock}
	\end{frame}
	
	\begin{frame}{Q123: Regulatory Inspection}
		\begin{block}{Question}
			A regulator requests documents during an inspection. What governance response is BEST?
		\end{block}
		\begin{enumerate}[(A)]
			\item Delay and minimize
			\item Respond fully and honestly per legal guidance
			\item Refuse
			\item Provide false documents
		\end{enumerate}
	\end{frame}
	
	\begin{frame}{Q123: Explanation}
		\begin{block}{Correct Answer: B}
			\textbf{Honest cooperation protects the enterprise:}
			\begin{itemize}
				\item Honest cooperation is required
				\item Legal guidance ensures accuracy
				\item Refusal and falsehood worsen exposure
			\end{itemize}
		\end{block}
		\begin{alertblock}{Critical Trap}
			Honest cooperation protects the enterprise.
		\end{alertblock}
	\end{frame}
	
	\begin{frame}{Q124: Contracting Ethics}
		\begin{block}{Question}
			A contract clause requires the organization to breach its own ethics code. What should governance do?
		\end{block}
		\begin{enumerate}[(A)]
			\item Comply silently
			\item Refuse, escalate, and renegotiate the clause
			\item Ignore the code
			\item Terminate immediately without discussion
		\end{enumerate}
	\end{frame}
	
	\begin{frame}{Q124: Explanation}
		\begin{block}{Correct Answer: B}
			\textbf{Ethics are non-negotiable:}
			\begin{itemize}
				\item Ethics clauses matter
				\item Renegotiation is the first step
				\item Silent breach damages integrity
			\end{itemize}
		\end{block}
		\begin{alertblock}{Critical Trap}
			Ethics are non-negotiable; contracts can be.
		\end{alertblock}
	\end{frame}
	
	\begin{frame}{Q125: Records Retention}
		\begin{block}{Question}
			A retention schedule conflicts with a legal hold. What should governance require?
		\end{block}
		\begin{enumerate}[(A)]
			\item Destroy records per schedule
			\item Preserve records under legal hold until released
			\item Ignore the hold
			\item Delete everything
		\end{enumerate}
	\end{frame}
	
	\begin{frame}{Q125: Explanation}
		\begin{block}{Correct Answer: B}
			\textbf{Legal holds override retention rules:}
			\begin{itemize}
				\item Legal holds override retention schedules
				\item Destruction during a hold is a serious violation
				\item Preservation protects the enterprise
			\end{itemize}
		\end{block}
		\begin{alertblock}{Critical Trap}
			Legal holds override retention rules.
		\end{alertblock}
	\end{frame}
	
	% ============================================================
	\section{3. Roles and Responsibilities}
	% ============================================================
	
	\begin{frame}{Q126: Risk Owner Authority}
		\begin{block}{Question}
			A risk owner lacks the authority to approve the required budget. What should happen?
		\end{block}
		\begin{enumerate}[(A)]
			\item Approve without authority
			\item Reassign ownership to someone with authority or escalate for authorization
			\item Ignore the risk
			\item Lower the rating
		\end{enumerate}
	\end{frame}
	
	\begin{frame}{Q126: Explanation}
		\begin{block}{Correct Answer: B}
			\textbf{Ownership needs authority to act:}
			\begin{itemize}
				\item Ownership requires authority
				\item Reassignment or escalation resolves the gap
				\item Lowering ratings hides the issue
			\end{itemize}
		\end{block}
		\begin{alertblock}{Critical Trap}
			Ownership needs authority to act.
		\end{alertblock}
	\end{frame}
	
	\begin{frame}{Q127: Data Ownership Handover}
		\begin{block}{Question}
			The data owner leaves the company. What is the MOST important governance action?
		\end{block}
		\begin{enumerate}[(A)]
			\item Delete the data
			\item Formally assign a new data owner with documented classification and access rights
			\item Ignore the departure
			\item Delegate to the vendor
		\end{enumerate}
	\end{frame}
	
	\begin{frame}{Q127: Explanation}
		\begin{block}{Correct Answer: B}
			\textbf{Data ownership must be handed over formally:}
			\begin{itemize}
				\item Formal handover preserves accountability
				\item Access rights must be updated
				\item Orphaned data risks exposure
			\end{itemize}
		\end{block}
		\begin{alertblock}{Critical Trap}
			Data ownership must be handed over formally.
		\end{alertblock}
	\end{frame}
	
	\begin{frame}{Q128: Three Lines in Crisis}
		\begin{block}{Question}
			During an incident, the first line is overwhelmed. Can the second line take over operational response?
		\end{block}
		\begin{enumerate}[(A)]
			\item Yes, without conditions
			\item Yes, if done through defined crisis procedures with clear authority
			\item No, never
			\item Only the third line may respond
		\end{enumerate}
	\end{frame}
	
	\begin{frame}{Q128: Explanation}
		\begin{block}{Correct Answer: B}
			\textbf{Crisis roles can shift with governance:}
			\begin{itemize}
				\item Crisis roles can shift if defined
				\item Clear authority prevents confusion
				\item Independence is preserved via process
			\end{itemize}
		\end{block}
		\begin{alertblock}{Critical Trap}
			Crisis roles can shift, but with governance.
		\end{alertblock}
	\end{frame}
	
	\begin{frame}{Q129: Audit Committee Independence}
		\begin{block}{Question}
			The audit committee chair also runs operations. What is the concern?
		\end{block}
		\begin{enumerate}[(A)]
			\item Stronger oversight
			\item Compromised independence and potential conflict of interest
			\item Lower cost
			\item Better reporting
		\end{enumerate}
	\end{frame}
	
	\begin{frame}{Q129: Explanation}
		\begin{block}{Correct Answer: B}
			\textbf{Independence cannot be compromised:}
			\begin{itemize}
				\item Independence is essential to oversight
				\item Combining operations and audit roles weakens challenge
				\item Governance requires separation
			\end{itemize}
		\end{block}
		\begin{alertblock}{Critical Trap}
			Independence cannot be compromised.
		\end{alertblock}
	\end{frame}
	
	\begin{frame}{Q130: Board Technology Advisor}
		\begin{block}{Question}
			The board hires an external technology advisor. What governance benefit is MOST likely?
		\end{block}
		\begin{enumerate}[(A)]
			\item Cheaper decisions
			\item Independent expertise to improve oversight quality
			\item Faster meetings
			\item Vendor preference
		\end{enumerate}
	\end{frame}
	
	\begin{frame}{Q130: Explanation}
		\begin{block}{Correct Answer: B}
			\textbf{Advisory expertise strengthens oversight:}
			\begin{itemize}
				\item Expertise deepens oversight
				\item Independence supports objectivity
				\item Cheaper is not the goal
			\end{itemize}
		\end{block}
		\begin{alertblock}{Critical Trap}
			Advisory expertise strengthens oversight.
		\end{alertblock}
	\end{frame}
	
	\begin{frame}{Q131: Executive Reporting Line}
		\begin{block}{Question}
			The CISO reports to the CIO, and the CIO is under investigation for control failures. What governance issue is MOST likely?
		\end{block}
		\begin{enumerate}[(A)]
			\item Independence of the CISO may be compromised
			\item Faster reporting
			\item Lower cost
			\item Better culture
		\end{enumerate}
	\end{frame}
	
	\begin{frame}{Q131: Explanation}
		\begin{block}{Correct Answer: A}
			\textbf{Reporting lines affect independence:}
			\begin{itemize}
				\item Independence protects security oversight
				\item Conflict may deter honest reporting
				\item Governance must review reporting lines
			\end{itemize}
		\end{block}
		\begin{alertblock}{Critical Trap}
			Reporting lines affect independence.
		\end{alertblock}
	\end{frame}
	
	\begin{frame}{Q132: Committee Quorum}
		\begin{block}{Question}
			A critical risk committee often fails to meet quorum. What governance action is BEST?
		\end{block}
		\begin{enumerate}[(A)]
			\item Ignore the failures
			\item Review membership, delegation, and meeting cadence to ensure quorum
			\item Cancel the committee
			\item Add more members without structure
		\end{enumerate}
	\end{frame}
	
	\begin{frame}{Q132: Explanation}
		\begin{block}{Correct Answer: B}
			\textbf{Committees must function:}
			\begin{itemize}
				\item Quorum ensures decisions
				\item Membership and delegation support attendance
				\item Adding members without structure may worsen
			\end{itemize}
		\end{block}
		\begin{alertblock}{Critical Trap}
			Committees must function.
		\end{alertblock}
	\end{frame}
	
	\begin{frame}{Q133: Delegated Authority}
		\begin{block}{Question}
			A manager approves expenditures above their delegated authority. What is the governance finding?
		\end{block}
		\begin{enumerate}[(A)]
			\item Acceptable
			\item A breach of delegated authority requiring review and corrective action
			\item Faster decisions
			\item Good practice
		\end{enumerate}
	\end{frame}
	
	\begin{frame}{Q133: Explanation}
		\begin{block}{Correct Answer: B}
			\textbf{Authority limits exist for a reason:}
			\begin{itemize}
				\item Limits define authority
				\item Breaches require investigation
				\item Tolerance for overreach erodes governance
			\end{itemize}
		\end{block}
		\begin{alertblock}{Critical Trap}
			Authority limits exist for a reason.
		\end{alertblock}
	\end{frame}
	
	\begin{frame}{Q134: Shadow Ownership}
		\begin{block}{Question}
			Multiple teams claim to own the same data. What is the MAIN governance risk?
		\end{block}
		\begin{enumerate}[(A)]
			\item Accountability confusion
			\item Faster processing
			\item Lower cost
			\item Better controls
		\end{enumerate}
	\end{frame}
	
	\begin{frame}{Q134: Explanation}
		\begin{block}{Correct Answer: A}
			\textbf{Single owner prevents gaps:}
			\begin{itemize}
				\item Ownership must be single and clear
				\item Confusion leads to gaps
				\item Controls degrade without ownership
			\end{itemize}
		\end{block}
		\begin{alertblock}{Critical Trap}
			Single owner prevents gaps.
		\end{alertblock}
	\end{frame}
	
	\begin{frame}{Q135: Governance of Shared Services}
		\begin{block}{Question}
			A shared service is used by many business units. Who is accountable for its risk?
		\end{block}
		\begin{enumerate}[(A)]
			\item The business unit with the largest budget
			\item A designated owner with cross-functional authority and clear accountability
			\item The vendor
			\item Nobody
		\end{enumerate}
	\end{frame}
	
	\begin{frame}{Q135: Explanation}
		\begin{block}{Correct Answer: B}
			\textbf{Shared services need single ownership:}
			\begin{itemize}
				\item Shared services need a single owner
				\item Cross-functional authority enables decisions
				\item Unclear ownership creates unmanaged risk
			\end{itemize}
		\end{block}
		\begin{alertblock}{Critical Trap}
			Shared services need single ownership.
		\end{alertblock}
	\end{frame}
	
	% ============================================================
	\section{4. Enterprise Risk Management}
	% ============================================================
	
	\begin{frame}{Q136: Risk Universe Refresh}
		\begin{block}{Question}
			The business enters a new market. What is the FIRST thing to update in the risk universe?
		\end{block}
		\begin{enumerate}[(A)]
			\item Vendor list
			\item Business objectives and dependencies introduced by the market entry
			\item Marketing plan
			\item Holiday schedule
		\end{enumerate}
	\end{frame}
	
	\begin{frame}{Q136: Explanation}
		\begin{block}{Correct Answer: B}
			\textbf{New markets reshape the risk universe:}
			\begin{itemize}
				\item Strategy changes the risk universe
				\item New dependencies emerge
				\item Objectives define priorities
			\end{itemize}
		\end{block}
		\begin{alertblock}{Critical Trap}
			New markets reshape the risk universe.
		\end{alertblock}
	\end{frame}
	
	\begin{frame}{Q137: Risk Aggregation}
		\begin{block}{Question}
			Aggregation across business units may hide a concentration. What is the BEST governance response?
		\end{block}
		\begin{enumerate}[(A)]
			\item Ignore the concentration
			\item Present both aggregate and drill-down views to reveal concentration
			\item Reduce reporting
			\item Increase budgets
		\end{enumerate}
	\end{frame}
	
	\begin{frame}{Q137: Explanation}
		\begin{block}{Correct Answer: B}
			\textbf{Aggregation must allow drill-down:}
			\begin{itemize}
				\item Aggregate views hide concentration
				\item Drill-down reveals exposure
				\item Governance needs visibility
			\end{itemize}
		\end{block}
		\begin{alertblock}{Critical Trap}
			Aggregation must allow drill-down.
		\end{alertblock}
	\end{frame}
	
	\begin{frame}{Q138: Emerging Risk Register}
		\begin{block}{Question}
			Which source BEST feeds the emerging risk register?
		\end{block}
		\begin{enumerate}[(A)]
			\item Historical incidents only
			\item Threat intelligence, regulatory horizon scanning, and strategic changes
			\item Vendor brochures
			\item Employee surveys only
		\end{enumerate}
	\end{frame}
	
	\begin{frame}{Q138: Explanation}
		\begin{block}{Correct Answer: B}
			\textbf{Emerging risks need forward-looking sources:}
			\begin{itemize}
				\item Emerging risks have little history
				\item Horizon scanning and intelligence reveal trends
				\item Brochures and surveys add context but not emerging risk signals
			\end{itemize}
		\end{block}
		\begin{alertblock}{Critical Trap}
			Emerging risks need forward-looking sources.
		\end{alertblock}
	\end{frame}
	
	\begin{frame}{Q139: Risk Tolerance Translation}
		\begin{block}{Question}
			The board tolerance is expressed in financial loss. How should this translate into operational metrics?
		\end{block}
		\begin{enumerate}[(A)]
			\item No translation needed
			\item Define operational indicators that correlate with financial tolerance
			\item Ignore the financial tolerance
			\item Use industry averages only
		\end{enumerate}
	\end{frame}
	
	\begin{frame}{Q139: Explanation}
		\begin{block}{Correct Answer: B}
			\textbf{Tolerance must become measurable:}
			\begin{itemize}
				\item Operational metrics must connect to financial limits
				\item Correlation supports early warning
				\item Ignoring tolerance breaks governance
			\end{itemize}
		\end{block}
		\begin{alertblock}{Critical Trap}
			Tolerance must become measurable operations.
		\end{alertblock}
	\end{frame}
	
	\begin{frame}{Q140: Risk Committee Decisions}
		\begin{block}{Question}
			A risk committee decision conflicts with a business unit's plan. What is the BEST path?
		\end{block}
		\begin{enumerate}[(A)]
			\item The business unit proceeds
			\item Escalate to the appropriate authority for resolution
			\item Ignore the committee
			\item Cancel the project
		\end{enumerate}
	\end{frame}
	
	\begin{frame}{Q140: Explanation}
		\begin{block}{Correct Answer: B}
			\textbf{Conflicts go to authority:}
			\begin{itemize}
				\item Conflict requires escalation
				\item Authority resolves conflicts
				\item Silent proceeding breaks governance
			\end{itemize}
		\end{block}
		\begin{alertblock}{Critical Trap}
			Conflicts go to authority, not silent override.
		\end{alertblock}
	\end{frame}
	
	\begin{frame}{Q141: Risk Profile Update}
		\begin{block}{Question}
			Risk profile updates should be triggered by:
		\end{block}
		\begin{enumerate}[(A)]
			\item Only annual cycles
			\item Material changes, incidents, and periodic cadence
			\item Only audits
			\item Only board meetings
		\end{enumerate}
	\end{frame}
	
	\begin{frame}{Q141: Explanation}
		\begin{block}{Correct Answer: B}
			\textbf{Update profiles on change:}
			\begin{itemize}
				\item Material changes alter the profile
				\item Incidents reveal gaps
				\item Periodic cadence prevents drift
			\end{itemize}
		\end{block}
		\begin{alertblock}{Critical Trap}
			Update profiles on change, not just on schedule.
		\end{alertblock}
	\end{frame}
	
	\begin{frame}{Q142: Risk Appetite Cascade}
		\begin{block}{Question}
			A board appetite statement exists but business units operate to different thresholds. What governance gap is MOST likely?
		\end{block}
		\begin{enumerate}[(A)]
			\item No gap
			\item Missing translation and communication of appetite to the operational level
			\item Too many thresholds
			\item Too few metrics
		\end{enumerate}
	\end{frame}
	
	\begin{frame}{Q142: Explanation}
		\begin{block}{Correct Answer: B}
			\textbf{Untranslated appetite drifts:}
			\begin{itemize}
				\item Appetite must be operationalized
				\item Without translation, units drift
				\item Communication aligns behavior
			\end{itemize}
		\end{block}
		\begin{alertblock}{Critical Trap}
			Untranslated appetite drifts at the operational level.
		\end{alertblock}
	\end{frame}
	
	\begin{frame}{Q143: Risk Reporting to Regulators}
		\begin{block}{Question}
			A regulator asks for evidence of risk oversight. What content is MOST valuable?
		\end{block}
		\begin{enumerate}[(A)]
			\item Marketing summary
			\item Evidence of governance structures, decisions, escalations, and remediations
			\item Vendor brochures
			\item Employee birthdays
		\end{enumerate}
	\end{frame}
	
	\begin{frame}{Q143: Explanation}
		\begin{block}{Correct Answer: B}
			\textbf{Oversight requires evidence:}
			\begin{itemize}
				\item Regulators need evidence of governance
				\item Decisions and escalations show oversight
				\item Remediations show action
			\end{itemize}
		\end{block}
		\begin{alertblock}{Critical Trap}
			Oversight requires evidence, not slogans.
		\end{alertblock}
	\end{frame}
	
	\begin{frame}{Q144: Risk Communication Quality}
		\begin{block}{Question}
			Which risk communication is MOST likely to trigger appropriate action?
		\end{block}
		\begin{enumerate}[(A)]
			\item Vague fear
			\item Clear risk, business consequence, options, and recommendation
			\item Technical jargon
			\item Too many details
		\end{enumerate}
	\end{frame}
	
	\begin{frame}{Q144: Explanation}
		\begin{block}{Correct Answer: B}
			\textbf{Clarity moves leaders to act:}
			\begin{itemize}
				\item Clear options drive action
				\item Business consequence aligns with decisions
				\item Jargon and volume obscure
			\end{itemize}
		\end{block}
		\begin{alertblock}{Critical Trap}
			Clarity moves leaders to act.
		\end{alertblock}
	\end{frame}
	
	\begin{frame}{Q145: Key Risk Indicators}
		\begin{block}{Question}
			A KRI has been green for years but a breach occurred. What is the MOST likely cause?
		\end{block}
		\begin{enumerate}[(A)]
			\item The KRI is too sensitive
			\item The KRI is not aligned to the actual risk driver or thresholds are too loose
			\item The KRI is too new
			\item The KRI is missing data
		\end{enumerate}
	\end{frame}
	
	\begin{frame}{Q145: Explanation}
		\begin{block}{Correct Answer: B}
			\textbf{A green KRI after a breach signals misalignment:}
			\begin{itemize}
				\item Misalignment hides exposure
				\item Loose thresholds miss signal
				\item KRIs must reflect real risk drivers
			\end{itemize}
		\end{block}
		\begin{alertblock}{Critical Trap}
			A green KRI after a breach signals misalignment.
		\end{alertblock}
	\end{frame}
	
	% ============================================================
	\section{5. Risk Culture}
	% ============================================================
	
	\begin{frame}{Q146: Cultural Silence}
		\begin{block}{Question}
			After a serious incident, no one reported the near-miss that preceded it. What culture failure is MOST likely?
		\end{block}
		\begin{enumerate}[(A)]
			\item Strong culture
			\item Fear of blame prevented reporting
			\item Excellent training
			\item Strong controls
		\end{enumerate}
	\end{frame}
	
	\begin{frame}{Q146: Explanation}
		\begin{block}{Correct Answer: B}
			\textbf{Fear silences early warnings:}
			\begin{itemize}
				\item Fear suppresses reporting
				\item Near-misses are early warnings
				\item Silence allows escalation into disasters
			\end{itemize}
		\end{block}
		\begin{alertblock}{Critical Trap}
			Fear silences early warnings.
		\end{alertblock}
	\end{frame}
	
	\begin{frame}{Q147: Positive Incentives}
		\begin{block}{Question}
			Which incentive BEST supports risk culture?
		\end{block}
		\begin{enumerate}[(A)]
			\item Rewarding only revenue
			\item Rewarding both strong results and risk-aware behavior
			\item Rewarding only cost reduction
			\item Rewarding silence
		\end{enumerate}
	\end{frame}
	
	\begin{frame}{Q147: Explanation}
		\begin{block}{Correct Answer: B}
			\textbf{Reward behavior you want:}
			\begin{itemize}
				\item Balanced rewards shape behavior
				\item Risk-aware behavior should be recognized
				\item One-sided rewards create blind spots
			\end{itemize}
		\end{block}
		\begin{alertblock}{Critical Trap}
			Reward behavior you want to see.
		\end{alertblock}
	\end{frame}
	
	\begin{frame}{Q148: Board Culture}
		\begin{block}{Question}
			The board rarely challenges executive reports. What culture risk is MOST likely?
		\end{block}
		\begin{enumerate}[(A)]
			\item Stronger governance
			\item Groupthink and reduced independent oversight
			\item Better training
			\item Faster decisions
		\end{enumerate}
	\end{frame}
	
	\begin{frame}{Q148: Explanation}
		\begin{block}{Correct Answer: B}
			\textbf{Boards must challenge:}
			\begin{itemize}
				\item Challenge is essential to oversight
				\item Silence allows drift
				\item Boards must ask tough questions
			\end{itemize}
		\end{block}
		\begin{alertblock}{Critical Trap}
			Boards must challenge, not just accept.
		\end{alertblock}
	\end{frame}
	
	\begin{frame}{Q149: Culture Measurement}
		\begin{block}{Question}
			Which is the MOST leading indicator of healthy risk culture?
		\end{block}
		\begin{enumerate}[(A)]
			\item Training completion rate
			\item Frequency of near-miss reporting
			\item Number of policy documents
			\item Vendor count
		\end{enumerate}
	\end{frame}
	
	\begin{frame}{Q149: Explanation}
		\begin{block}{Correct Answer: B}
			\textbf{Near-miss reporting measures trust:}
			\begin{itemize}
				\item Near-miss reporting signals trust
				\item Training is an input, not an outcome
				\item Document count does not indicate culture
			\end{itemize}
		\end{block}
		\begin{alertblock}{Critical Trap}
			Near-miss reporting measures trust.
		\end{alertblock}
	\end{frame}
	
	\begin{frame}{Q150: Executive Exception}
		\begin{block}{Question}
			An executive routinely bypasses control processes. What culture impact is MOST likely?
		\end{block}
		\begin{enumerate}[(A)]
			\item Stronger controls
			\item Others imitate the exception, weakening controls
			\item Faster projects
			\item Better compliance
		\end{enumerate}
	\end{frame}
	
	\begin{frame}{Q150: Explanation}
		\begin{block}{Correct Answer: B}
			\textbf{Executive exceptions cascade downward:}
			\begin{itemize}
				\item Leaders set the norm
				\item Exceptions normalize bypasses
				\item Controls decay from the top
			\end{itemize}
		\end{block}
		\begin{alertblock}{Critical Trap}
			Executive exceptions cascade downward.
		\end{alertblock}
	\end{frame}
	
	\begin{frame}{Q151: Awareness Resistance}
		\begin{block}{Question}
			Employees resist security awareness training. What should governance do FIRST?
		\end{block}
		\begin{enumerate}[(A)]
			\item Mandate and punish
			\item Understand the resistance and tailor the program to roles and relevance
			\item Ignore
			\item Reduce training
		\end{enumerate}
	\end{frame}
	
	\begin{frame}{Q151: Explanation}
		\begin{block}{Correct Answer: B}
			\textbf{Tailor training:}
			\begin{itemize}
				\item Understanding improves engagement
				\item Tailoring makes training relevant
				\item Punishment alone breeds resentment
			\end{itemize}
		\end{block}
		\begin{alertblock}{Critical Trap}
			Tailor training; do not just mandate.
		\end{alertblock}
	\end{frame}
	
	\begin{frame}{Q152: Risk Communication}
		\begin{block}{Question}
			An organization uses fear tactics for risk communication. What is the likely long-term impact?
		\end{block}
		\begin{enumerate}[(A)]
			\item Stronger trust
			\item Desensitization and reduced engagement
			\item Faster action
			\item Better culture
		\end{enumerate}
	\end{frame}
	
	\begin{frame}{Q152: Explanation}
		\begin{block}{Correct Answer: B}
			\textbf{Fear fatigues:}
			\begin{itemize}
				\item Fear loses effect
				\item Engagement drops
				\item Trust erodes
			\end{itemize}
		\end{block}
		\begin{alertblock}{Critical Trap}
			Fear fatigues; engagement sustains.
		\end{alertblock}
	\end{frame}
	
	\begin{frame}{Q153: Culture After Incident}
		\begin{block}{Question}
			After an incident, leaders focus on blame. What culture risk is HIGHEST?
		\end{block}
		\begin{enumerate}[(A)]
			\item Stronger reporting
			\item Reduced future reporting and learning
			\item Faster recovery
			\item Better controls
		\end{enumerate}
	\end{frame}
	
	\begin{frame}{Q153: Explanation}
		\begin{block}{Correct Answer: B}
			\textbf{Blame blocks learning:}
			\begin{itemize}
				\item Blame suppresses disclosure
				\item Learning requires safety
				\item Reporting shrinks
			\end{itemize}
		\end{block}
		\begin{alertblock}{Critical Trap}
			Blame blocks learning.
		\end{alertblock}
	\end{frame}
	
	\begin{frame}{Q154: Ethics Training}
		\begin{block}{Question}
			Ethics training exists but is a checkbox exercise. What is the MAIN weakness?
		\end{block}
		\begin{enumerate}[(A)]
			\item Too long
			\item It does not change behavior or reinforce culture
			\item Too few sessions
			\item Wrong location
		\end{enumerate}
	\end{frame}
	
	\begin{frame}{Q154: Explanation}
		\begin{block}{Correct Answer: B}
			\textbf{Checkbox ethics is not ethics:}
			\begin{itemize}
				\item Culture requires behavior
				\item Checkbox training does not reinforce values
				\item Interactive and scenario-based training is stronger
			\end{itemize}
		\end{block}
		\begin{alertblock}{Critical Trap}
			Checkbox ethics is not ethics.
		\end{alertblock}
	\end{frame}
	
	\begin{frame}{Q155: Leadership Modeling}
		\begin{block}{Question}
			Leaders participate in risk training. What is the MOST likely culture benefit?
		\end{block}
		\begin{enumerate}[(A)]
			\item Lower cost
			\item Increased credibility and shared ownership of risk
			\item Faster projects
			\item More audits
		\end{enumerate}
	\end{frame}
	
	\begin{frame}{Q155: Explanation}
		\begin{block}{Correct Answer: B}
			\textbf{Leaders modeling risk builds culture:}
			\begin{itemize}
				\item Leaders modeling behavior builds trust
				\item Shared ownership strengthens culture
				\item Participation signals priority
			\end{itemize}
		\end{block}
		\begin{alertblock}{Critical Trap}
			Leaders modeling risk builds culture.
		\end{alertblock}
	\end{frame}
	
	% ============================================================
	\section{6. Third-Party Governance}
	% ============================================================
	
	\begin{frame}{Q156: Tiering Criteria}
		\begin{block}{Question}
			Vendors are tiered by contract value. What is the MOST important governance flaw?
		\end{block}
		\begin{enumerate}[(A)]
			\item Too many tiers
			\item Contract value is not the same as risk; criticality and data sensitivity should drive tiering
			\item Too few tiers
			\item Too few vendors
		\end{enumerate}
	\end{frame}
	
	\begin{frame}{Q156: Explanation}
		\begin{block}{Correct Answer: B}
			\textbf{Tier by risk, not spend:}
			\begin{itemize}
				\item Risk drives tiering
				\item Contract value alone misleads
				\item Tiering should reflect exposure
			\end{itemize}
		\end{block}
		\begin{alertblock}{Critical Trap}
			Tier by risk, not spend.
		\end{alertblock}
	\end{frame}
	
	\begin{frame}{Q157: Fourth-Party Visibility}
		\begin{block}{Question}
			A vendor cannot name its subcontractors. What is the MOST likely governance concern?
		\end{block}
		\begin{enumerate}[(A)]
			\item Hidden risk and weak oversight of the extended supply chain
			\item Too many subcontractors
			\item Too few invoices
			\item Too much training
		\end{enumerate}
	\end{frame}
	
	\begin{frame}{Q157: Explanation}
		\begin{block}{Correct Answer: A}
			\textbf{No visibility, no oversight:}
			\begin{itemize}
				\item Visibility enables oversight
				\item Unknown subcontractors create risk
				\item Governance requires disclosure
			\end{itemize}
		\end{block}
		\begin{alertblock}{Critical Trap}
			No visibility, no oversight.
		\end{alertblock}
	\end{frame}
	
	\begin{frame}{Q158: Vendor Exit Execution}
		\begin{block}{Question}
			A critical vendor relationship ends abruptly. What governance should have prevented the disruption?
		\end{block}
		\begin{enumerate}[(A)]
			\item A larger contract
			\item A pre-agreed exit plan with data return and continuity steps
			\item A faster payment
			\item More audits
		\end{enumerate}
	\end{frame}
	
	\begin{frame}{Q158: Explanation}
		\begin{block}{Correct Answer: B}
			\textbf{Exit plans are part of onboarding:}
			\begin{itemize}
				\item Exit plans reduce disruption
				\item Data return is essential
				\item Continuity protects operations
			\end{itemize}
		\end{block}
		\begin{alertblock}{Critical Trap}
			Exit plans are part of onboarding.
		\end{alertblock}
	\end{frame}
	
	\begin{frame}{Q159: Right to Audit}
		\begin{block}{Question}
			A critical vendor's contract lacks a right-to-audit clause. What is the governance impact?
		\end{block}
		\begin{enumerate}[(A)]
			\item Stronger oversight
			\item Limited ability to verify vendor controls
			\item Lower risk
			\item Faster response
		\end{enumerate}
	\end{frame}
	
	\begin{frame}{Q159: Explanation}
		\begin{block}{Correct Answer: B}
			\textbf{No audit rights, no proof:}
			\begin{itemize}
				\item Verification depends on access
				\item Audits provide evidence
				\item Without rights, oversight is limited
			\end{itemize}
		\end{block}
		\begin{alertblock}{Critical Trap}
			No audit rights, no proof.
		\end{alertblock}
	\end{frame}
	
	\begin{frame}{Q160: Vendor Certification}
		\begin{block}{Question}
			A vendor holds a certification but recently suffered a breach. What governance lesson applies?
		\end{block}
		\begin{enumerate}[(A)]
			\item Certifications guarantee safety
			\item Certifications indicate capability at a point in time; ongoing oversight remains necessary
			\item Breaches are impossible with certification
			\item Certifications should be abandoned
		\end{enumerate}
	\end{frame}
	
	\begin{frame}{Q160: Explanation}
		\begin{block}{Correct Answer: B}
			\textbf{Certification is not a guarantee:}
			\begin{itemize}
				\item Certifications are snapshots
				\item Continuous oversight is required
				\item Breaches can still occur
			\end{itemize}
		\end{block}
		\begin{alertblock}{Critical Trap}
			Certification is not a guarantee.
		\end{alertblock}
	\end{frame}
	
	\begin{frame}{Q161: Vendor Notifications}
		\begin{block}{Question}
			The vendor contract allows 30 days to notify of a breach. What governance concern is most relevant?
		\end{block}
		\begin{enumerate}[(A)]
			\item Faster response
			\item Delayed notification limits the organization's ability to act quickly
			\item Better compliance
			\item Lower cost
		\end{enumerate}
	\end{frame}
	
	\begin{frame}{Q161: Explanation}
		\begin{block}{Correct Answer: B}
			\textbf{Slow notification reduces response time:}
			\begin{itemize}
				\item Timely notification enables response
				\item Long windows increase exposure
				\item Contracts should align to regulatory timelines
			\end{itemize}
		\end{block}
		\begin{alertblock}{Critical Trap}
			Slow notification reduces response time.
		\end{alertblock}
	\end{frame}
	
	\begin{frame}{Q162: Vendor Concentration}
		\begin{block}{Question}
			Three critical vendors depend on the same cloud provider. What risk should be assessed FIRST?
		\end{block}
		\begin{enumerate}[(A)]
			\item Marketing strategy
			\item Concentration risk at the underlying provider level
			\item HR planning
			\item Catering
		\end{enumerate}
	\end{frame}
	
	\begin{frame}{Q162: Explanation}
		\begin{block}{Correct Answer: B}
			\textbf{Look beyond direct vendors:}
			\begin{itemize}
				\item Underlying concentration amplifies impact
				\item Shared dependency increases single points of failure
				\item Governance must assess chain-wide risk
			\end{itemize}
		\end{block}
		\begin{alertblock}{Critical Trap}
			Look beyond direct vendors to underlying dependencies.
		\end{alertblock}
	\end{frame}
	
	\begin{frame}{Q163: Vendor Governance Cadence}
		\begin{block}{Question}
			Vendor reviews are quarterly for critical vendors and annual for others. Is this governance sound?
		\end{block}
		\begin{enumerate}[(A)]
			\item No, all should be annual
			\item Yes, tiered cadence aligned to risk is appropriate
			\item No, all should be quarterly
			\item No, no reviews needed
		\end{enumerate}
	\end{frame}
	
	\begin{frame}{Q163: Explanation}
		\begin{block}{Correct Answer: B}
			\textbf{Cadence follows criticality:}
			\begin{itemize}
				\item Cadence should match criticality
				\item Critical vendors need closer oversight
				\item Risk-based review is efficient and effective
			\end{itemize}
		\end{block}
		\begin{alertblock}{Critical Trap}
			Cadence follows criticality.
		\end{alertblock}
	\end{frame}
	
	\begin{frame}{Q164: Vendor Incentives}
		\begin{block}{Question}
			A vendor contract pays more for faster delivery only. What governance risk is created?
		\end{block}
		\begin{enumerate}[(A)]
			\item Faster delivery is always safe
			\item Incentives may encourage cutting corners on security
			\item Fewer suppliers
			\item Higher costs
		\end{enumerate}
	\end{frame}
	
	\begin{frame}{Q164: Explanation}
		\begin{block}{Correct Answer: B}
			\textbf{Incentives shape vendor behavior:}
			\begin{itemize}
				\item One-sided incentives skew behavior
				\item Security must be rewarded too
				\item Balanced incentives protect the enterprise
			\end{itemize}
		\end{block}
		\begin{alertblock}{Critical Trap}
			Incentives shape vendor behavior.
		\end{alertblock}
	\end{frame}
	
	\begin{frame}{Q165: Vendor Data Return}
		\begin{block}{Question}
			A vendor is required to return all data at contract termination. What governance mechanism BEST verifies this?
		\end{block}
		\begin{enumerate}[(A)]
			\item Trust the vendor
			\item Certificate of destruction and audit rights to verify
			\item Ignore
			\item Delay return
		\end{enumerate}
	\end{frame}
	
	\begin{frame}{Q165: Explanation}
		\begin{block}{Correct Answer: B}
			\textbf{Trust, but verify data return:}
			\begin{itemize}
				\item Verification is essential
				\item Certificates provide evidence
				\item Audit rights support oversight
			\end{itemize}
		\end{block}
		\begin{alertblock}{Critical Trap}
			Trust, but verify data return.
		\end{alertblock}
	\end{frame}
	
	% ============================================================
	\section{7. Risk Appetite, Tolerance and Reporting}
	% ============================================================
	
	\begin{frame}{Q166: Board Reporting Escalation}
		\begin{block}{Question}
			The enterprise faces a material risk that has not been reported to the board. What is the MOST likely governance failure?
		\end{block}
		\begin{enumerate}[(A)]
			\item Too many dashboards
			\item Missing escalation and reporting paths for material risks
			\item Too few vendors
			\item Too much training
		\end{enumerate}
	\end{frame}
	
	\begin{frame}{Q166: Explanation}
		\begin{block}{Correct Answer: B}
			\textbf{Material risk belongs on the board's desk:}
			\begin{itemize}
				\item Material risk must reach the board
				\item Escalation paths ensure visibility
				\item Silence is a governance failure
			\end{itemize}
		\end{block}
		\begin{alertblock}{Critical Trap}
			Material risk belongs on the board's desk.
		\end{alertblock}
	\end{frame}
	
	\begin{frame}{Q167: Appetite for Innovation}
		\begin{block}{Question}
			A board sets a low appetite for operational risk but a high appetite for innovation. What governance guidance follows?
		\end{block}
		\begin{enumerate}[(A)]
			\item Avoid innovation
			\item Encourage innovation within defined risk thresholds and governance gates
			\item Ignore risk
			\item Cancel new ideas
		\end{enumerate}
	\end{frame}
	
	\begin{frame}{Q167: Explanation}
		\begin{block}{Correct Answer: B}
			\textbf{Balance innovation with risk appetite:}
			\begin{itemize}
				\item Appetite can differ by risk type
				\item Gates and thresholds shape safe innovation
				\item Avoidance and ignoring risk both fail
			\end{itemize}
		\end{block}
		\begin{alertblock}{Critical Trap}
			Balance innovation with risk appetite.
		\end{alertblock}
	\end{frame}
	
	\begin{frame}{Q168: Tolerance Breach}
		\begin{block}{Question}
			A KRI breaches tolerance. What is the FIRST governance expectation?
		\end{block}
		\begin{enumerate}[(A)]
			\item No action
			\item Defined escalation and response per governance
			\item New KRIs
			\item An audit
		\end{enumerate}
	\end{frame}
	
	\begin{frame}{Q168: Explanation}
		\begin{block}{Correct Answer: B}
			\textbf{Tolerance breaches need response:}
			\begin{itemize}
				\item Breaches should trigger action
				\item Defined responses prevent confusion
				\item Governance requires discipline
			\end{itemize}
		\end{block}
		\begin{alertblock}{Critical Trap}
			Tolerance breaches need response.
		\end{alertblock}
	\end{frame}
	
	\begin{frame}{Q169: Reporting Insight}
		\begin{block}{Question}
			The board receives 300 KRI values. What should governance change FIRST?
		\end{block}
		\begin{enumerate}[(A)]
			\item Add more KRIs
			\item Focus on a small set of key indicators with clear ownership and thresholds
			\item Reduce the reporting cadence
			\item Increase vendor count
		\end{enumerate}
	\end{frame}
	
	\begin{frame}{Q169: Explanation}
		\begin{block}{Correct Answer: B}
			\textbf{Fewer, better indicators:}
			\begin{itemize}
				\item Focused indicators drive decisions
				\item Ownership clarifies action
				\item Volume obscures signal
			\end{itemize}
		\end{block}
		\begin{alertblock}{Critical Trap}
			Fewer, better indicators.
		\end{alertblock}
	\end{frame}
	
	\begin{frame}{Q170: Loss Trending}
		\begin{block}{Question}
			A risk has been stable for five years but conditions are changing. What should governance require?
		\end{block}
		\begin{enumerate}[(A)]
			\item No action
			\item Reassessment and updated indicators
			\item Continue as usual
			\item Increase ratings
		\end{enumerate}
	\end{frame}
	
	\begin{frame}{Q170: Explanation}
		\begin{block}{Correct Answer: B}
			\textbf{Stable past is not a stable future:}
			\begin{itemize}
				\item Stable history may not reflect current risk
				\item Reassessment catches changes
				\item Adjustments follow data
			\end{itemize}
		\end{block}
		\begin{alertblock}{Critical Trap}
			Stable past is not a stable future.
		\end{alertblock}
	\end{frame}
	
	\begin{frame}{Q171: Assurance Source}
		\begin{block}{Question}
			A board wants confidence in a critical system's security. Which assurance source is MOST credible?
		\end{block}
		\begin{enumerate}[(A)]
			\item A vendor's marketing statement
			\item Independent assurance with evidence
			\item The system owner's assurance only
			\item An informal email
		\end{enumerate}
	\end{frame}
	
	\begin{frame}{Q171: Explanation}
		\begin{block}{Correct Answer: B}
			\textbf{Evidence and independence build trust:}
			\begin{itemize}
				\item Independence and evidence support credibility
				\item Marketing and informal statements lack rigor
				\item Self-assessment alone may be biased
			\end{itemize}
		\end{block}
		\begin{alertblock}{Critical Trap}
			Evidence and independence build trust.
		\end{alertblock}
	\end{frame}
	
	\begin{frame}{Q172: Appetite Alignment}
		\begin{block}{Question}
			Business unit investments are expanding beyond appetite. What governance action is MOST appropriate?
		\end{block}
		\begin{enumerate}[(A)]
			\item Ignore the expansion
			\item Re-affirm or adjust appetite and align investments to it
			\item Increase approvals
			\item Cancel all expansion
		\end{enumerate}
	\end{frame}
	
	\begin{frame}{Q172: Explanation}
		\begin{block}{Correct Answer: B}
			\textbf{Investments align with appetite:}
			\begin{itemize}
				\item Appetite must guide investments
				\item Alignment is a leadership decision
				\item Cancellation needs analysis
			\end{itemize}
		\end{block}
		\begin{alertblock}{Critical Trap}
			Investments align with appetite, not the reverse.
		\end{alertblock}
	\end{frame}
	
	\begin{frame}{Q173: Risk Statement Clarity}
		\begin{block}{Question}
			Which risk statement is MOST likely to prompt action from executives?
		\end{block}
		\begin{enumerate}[(A)]
			\item Some IT risks exist
			\item A breach could expose customer records, trigger regulatory fines, and damage brand trust within six months
			\item Firewall rules are outdated
			\item Log data is large
		\end{enumerate}
	\end{frame}
	
	\begin{frame}{Q173: Explanation}
		\begin{block}{Correct Answer: B}
			\textbf{Consequence and timing drive action:}
			\begin{itemize}
				\item Business consequence drives urgency
				\item Timeframe and impact create meaning
				\item Technical detail and vagueness lack pull
			\end{itemize}
		\end{block}
		\begin{alertblock}{Critical Trap}
			Consequence and timing drive action.
		\end{alertblock}
	\end{frame}
	
	\begin{frame}{Q174: Aggregation Risk}
		\begin{block}{Question}
			Combined risk exceeds appetite though each unit is within appetite. What does this reveal?
		\end{block}
		\begin{enumerate}[(A)]
			\item Appetite is too strict
			\item Aggregation reveals exposure that individual views miss
			\item The units are wrong
			\item No issue
		\end{enumerate}
	\end{frame}
	
	\begin{frame}{Q174: Explanation}
		\begin{block}{Correct Answer: B}
			\textbf{Portfolio risk beats unit-by-unit views:}
			\begin{itemize}
				\item Combined risk can exceed individual limits
				\item Aggregation reveals systemic exposure
				\item Appetite should consider the portfolio
			\end{itemize}
		\end{block}
		\begin{alertblock}{Critical Trap}
			Portfolio risk beats unit-by-unit views.
		\end{alertblock}
	\end{frame}
	
	\begin{frame}{Q175: Board Challenge}
		\begin{block}{Question}
			The board asks how a rating was determined. What is the BEST governance response?
		\end{block}
		\begin{enumerate}[(A)]
			\item Provide the methodology, data, and assumptions
			\item Decline to answer
			\item Reissue the rating higher
			\item Postpone
		\end{enumerate}
	\end{frame}
	
	\begin{frame}{Q175: Explanation}
		\begin{block}{Correct Answer: A}
			\textbf{Explain the methodology:}
			\begin{itemize}
				\item Transparency builds trust
				\item Methodology and data support credibility
				\item Avoidance undermines confidence
			\end{itemize}
		\end{block}
		\begin{alertblock}{Critical Trap}
			Explain the methodology, not just the number.
		\end{alertblock}
	\end{frame}
	
	% ============================================================
	\section{8. Business Continuity Governance}
	% ============================================================
	
	\begin{frame}{Q176: BCP Ownership}
		\begin{block}{Question}
			The BCP is owned by IT alone. What governance issue is MOST likely?
		\end{block}
		\begin{enumerate}[(A)]
			\item Too much scope
			\item Business process priorities and impact may be missing
			\item Overfunding
			\item Too many tests
		\end{enumerate}
	\end{frame}
	
	\begin{frame}{Q176: Explanation}
		\begin{block}{Correct Answer: B}
			\textbf{BCP must be business-owned:}
			\begin{itemize}
				\item Continuity is business-owned
				\item IT supports recovery
				\item IT-only plans miss business impact
			\end{itemize}
		\end{block}
		\begin{alertblock}{Critical Trap}
			BCP must be business-owned.
		\end{alertblock}
	\end{frame}
	
	\begin{frame}{Q177: BCP Test Results}
		\begin{block}{Question}
			A BCP test succeeds but the scenario was minor. What governance concern remains?
		\end{block}
		\begin{enumerate}[(A)]
			\item No concern
			\item The plan may not hold under realistic or severe scenarios
			\item Too many tests
			\item Too much budget
		\end{enumerate}
	\end{frame}
	
	\begin{frame}{Q177: Explanation}
		\begin{block}{Correct Answer: B}
			\textbf{Test realistic and severe scenarios:}
			\begin{itemize}
				\item Test severity matters
				\item Realistic scenarios reveal weaknesses
				\item Minor tests can create false confidence
			\end{itemize}
		\end{block}
		\begin{alertblock}{Critical Trap}
			Test realistic and severe scenarios.
		\end{alertblock}
	\end{frame}
	
	\begin{frame}{Q178: Continuity Roles}
		\begin{block}{Question}
			During a crisis, nobody is sure who is in charge. What is the governance gap?
		\end{block}
		\begin{enumerate}[(A)]
			\item Too many tools
			\item Undefined crisis roles and authority
			\item Too many policies
			\item Too few vendors
		\end{enumerate}
	\end{frame}
	
	\begin{frame}{Q178: Explanation}
		\begin{block}{Correct Answer: B}
			\textbf{Crisis needs clear command:}
			\begin{itemize}
				\item Crisis needs clear roles
				\item Ambiguity slows response
				\item Authority protects decisions
			\end{itemize}
		\end{block}
		\begin{alertblock}{Critical Trap}
			Crisis needs clear command.
		\end{alertblock}
	\end{frame}
	
	\begin{frame}{Q179: Plan Maintenance}
		\begin{block}{Question}
			The BCP was last updated four years ago. What is governance's BEST action?
		\end{block}
		\begin{enumerate}[(A)]
			\item Approve it as-is
			\item Require a review and update to reflect current operations
			\item Cancel the plan
			\item Archive it
		\end{enumerate}
	\end{frame}
	
	\begin{frame}{Q179: Explanation}
		\begin{block}{Correct Answer: B}
			\textbf{Stale plans fail when needed:}
			\begin{itemize}
				\item Plans must reflect current operations
				\item Reviews keep plans relevant
				\item Archiving loses readiness
			\end{itemize}
		\end{block}
		\begin{alertblock}{Critical Trap}
			Stale plans fail when needed.
		\end{alertblock}
	\end{frame}
	
	\begin{frame}{Q180: Continuity Governance}
		\begin{block}{Question}
			The board approves the BCP but never reviews test results. What governance gap is MOST likely?
		\end{block}
		\begin{enumerate}[(A)]
			\item Too much oversight
			\item Oversight lacks evidence of readiness
			\item Too few plans
			\item Too many auditors
		\end{enumerate}
	\end{frame}
	
	\begin{frame}{Q180: Explanation}
		\begin{block}{Correct Answer: B}
			\textbf{Oversight requires evidence of readiness:}
			\begin{itemize}
				\item Oversight needs test evidence
				\item Reviews confirm readiness
				\item Silence implies assumption
			\end{itemize}
		\end{block}
		\begin{alertblock}{Critical Trap}
			Oversight requires evidence of readiness.
		\end{alertblock}
	\end{frame}
	
	% ============================================================
	\section{9. IT Resource and Performance Governance}
	% ============================================================
	
	\begin{frame}{Q181: Investment Oversight}
		\begin{block}{Question}
			A major IT investment is approved without a business case. What is the governance concern?
		\end{block}
		\begin{enumerate}[(A)]
			\item Speed
			\item Value and risk are not demonstrated before spending
			\item Fewer approvals
			\item Better culture
		\end{enumerate}
	\end{frame}
	
	\begin{frame}{Q181: Explanation}
		\begin{block}{Correct Answer: B}
			\textbf{No business case, no oversight:}
			\begin{itemize}
				\item Business cases support value
				\item Risk and return should be shown
				\item Approval without analysis is a governance gap
			\end{itemize}
		\end{block}
		\begin{alertblock}{Critical Trap}
			No business case, no proper oversight.
		\end{alertblock}
	\end{frame}
	
	\begin{frame}{Q182: Performance Metrics}
		\begin{block}{Question}
			IT performance is measured only on cost. What governance risk is created?
		\end{block}
		\begin{enumerate}[(A)]
			\item Lower value and higher long-term risk
			\item Stronger oversight
			\item Faster delivery
			\item Better controls
		\end{enumerate}
	\end{frame}
	
	\begin{frame}{Q182: Explanation}
		\begin{block}{Correct Answer: A}
			\textbf{Cost alone is misleading:}
			\begin{itemize}
				\item Cost-only metrics ignore value and risk
				\item Balanced measures protect outcomes
				\item Single-metric focus creates blind spots
			\end{itemize}
		\end{block}
		\begin{alertblock}{Critical Trap}
			Cost alone is a misleading metric.
		\end{alertblock}
	\end{frame}
	
	\begin{frame}{Q183: Change Governance}
		\begin{block}{Question}
			A production change occurs without change control. What is the MOST likely governance risk?
		\end{block}
		\begin{enumerate}[(A)]
			\item Faster delivery
			\item Unassessed risk and untested changes in production
			\item Lower cost
			\item Better compliance
		\end{enumerate}
	\end{frame}
	
	\begin{frame}{Q183: Explanation}
		\begin{block}{Correct Answer: B}
			\textbf{Change governance prevents surprises:}
			\begin{itemize}
				\item Changes need review and testing
				\item Bypassing control increases risk
				\item Speed without governance creates exposure
			\end{itemize}
		\end{block}
		\begin{alertblock}{Critical Trap}
			Change governance prevents surprises.
		\end{alertblock}
	\end{frame}
	
	\begin{frame}{Q184: Resource Allocation}
		\begin{block}{Question}
			Resource decisions favor short-term projects over foundational controls. What is the governance risk?
		\end{block}
		\begin{enumerate}[(A)]
			\item Faster short-term results but increased long-term exposure
			\item Stronger controls
			\item Better culture
			\item Lower risk
		\end{enumerate}
	\end{frame}
	
	\begin{frame}{Q184: Explanation}
		\begin{block}{Correct Answer: A}
			\textbf{Short-term wins can hide long-term risk:}
			\begin{itemize}
				\item Foundational work protects the enterprise
				\item Short-term focus can erode resilience
				\item Balanced investment matters
			\end{itemize}
		\end{block}
		\begin{alertblock}{Critical Trap}
			Short-term wins can hide long-term risk.
		\end{alertblock}
	\end{frame}
	
	\begin{frame}{Q185: Value Realization}
		\begin{block}{Question}
			A project is completed but benefits are not tracked. What governance improvement is BEST?
		\end{block}
		\begin{enumerate}[(A)]
			\item Approve more projects
			\item Require post-implementation benefits and risk reviews
			\item Cancel the project
			\item Increase the budget
		\end{enumerate}
	\end{frame}
	
	\begin{frame}{Q185: Explanation}
		\begin{block}{Correct Answer: B}
			\textbf{Completion is not the same as value:}
			\begin{itemize}
				\item Value must be realized
				\item Reviews confirm outcomes
				\item Silent completion may hide failure
			\end{itemize}
		\end{block}
		\begin{alertblock}{Critical Trap}
			Completion is not the same as value.
		\end{alertblock}
	\end{frame}
	
	% ============================================================
	\section{10. Advanced Governance Scenarios}
	% ============================================================
	
	\begin{frame}{Q186: Board Technology Committee}
		\begin{block}{Question}
			A technology committee exists but lacks decision authority. What governance improvement is BEST?
		\end{block}
		\begin{enumerate}[(A)]
			\item Increase meeting frequency
			\item Define authority, escalation paths, and reporting lines to the board
			\item Add more members
			\item Reduce meetings
		\end{enumerate}
	\end{frame}
	
	\begin{frame}{Q186: Explanation}
		\begin{block}{Correct Answer: B}
			\textbf{Committees need authority:}
			\begin{itemize}
				\item Authority drives action
				\item Escalation ensures visibility
				\item More meetings alone add no power
			\end{itemize}
		\end{block}
		\begin{alertblock}{Critical Trap}
			Committees need authority and escalation.
		\end{alertblock}
	\end{frame}
	
	\begin{frame}{Q187: Independent Advice}
		\begin{block}{Question}
			The board suspects management is downplaying a risk. What governance action is MOST appropriate?
		\end{block}
		\begin{enumerate}[(A)]
			\item Accept management's view
			\item Seek independent assurance
			\item Ignore the concern
			\item Cancel the project
		\end{enumerate}
	\end{frame}
	
	\begin{frame}{Q187: Explanation}
		\begin{block}{Correct Answer: B}
			\textbf{Independence tests management:}
			\begin{itemize}
				\item Independent assurance tests management's view
				\item Boards need objective evidence
				\item Ignoring concerns increases exposure
			\end{itemize}
		\end{block}
		\begin{alertblock}{Critical Trap}
			Independence tests what management says.
		\end{alertblock}
	\end{frame}
	
	\begin{frame}{Q188: Executive Conflict}
		\begin{block}{Question}
			Two executives disagree on risk acceptance. What should happen NEXT?
		\end{block}
		\begin{enumerate}[(A)]
			\item The more senior wins
			\item Escalate to a designated decision authority with documented rationale
			\item Ignore
			\item Cancel
		\end{enumerate}
	\end{frame}
	
	\begin{frame}{Q188: Explanation}
		\begin{block}{Correct Answer: B}
			\textbf{Conflicts need structured resolution:}
			\begin{itemize}
				\item Conflict needs a defined authority
				\item Documentation preserves accountability
				\item Silence leaves risk unresolved
			\end{itemize}
		\end{block}
		\begin{alertblock}{Critical Trap}
			Conflicts need structured resolution.
		\end{alertblock}
	\end{frame}
	
	\begin{frame}{Q189: Risk Reporting Integrity}
		\begin{block}{Question}
			A risk report is edited to remove negative findings. What governance issue is MOST visible?
		\end{block}
		\begin{enumerate}[(A)]
			\item Strong oversight
			\item Loss of reporting integrity and trust
			\item Faster action
			\item Lower costs
		\end{enumerate}
	\end{frame}
	
	\begin{frame}{Q189: Explanation}
		\begin{block}{Correct Answer: B}
			\textbf{Report integrity is governance integrity:}
			\begin{itemize}
				\item Editing hides exposure
				\item Boards need honest information
				\item Trust erodes quickly
			\end{itemize}
		\end{block}
		\begin{alertblock}{Critical Trap}
			Report integrity is governance integrity.
		\end{alertblock}
	\end{frame}
	
	\begin{frame}{Q190: Governance of Vendors}
		\begin{block}{Question}
			A vendor's role is expanded beyond the approved scope without governance approval. What is the governance gap?
		\end{block}
		\begin{enumerate}[(A)]
			\item Scope creep without approval
			\item Better contracts
			\item Faster delivery
			\item Strong culture
		\end{enumerate}
	\end{frame}
	
	\begin{frame}{Q190: Explanation}
		\begin{block}{Correct Answer: A}
			\textbf{Scope changes need approval:}
			\begin{itemize}
				\item Expanded scope changes exposure
				\item Approval ensures oversight
				\item Silent scope change breaks governance
			\end{itemize}
		\end{block}
		\begin{alertblock}{Critical Trap}
			Scope changes need approval.
		\end{alertblock}
	\end{frame}
	
	\begin{frame}{Q191: Crisis Decision Rights}
		\begin{block}{Question}
			During an incident, decisions are delayed by unclear authority. What should governance define?
		\end{block}
		\begin{enumerate}[(A)]
			\item More meetings
			\item Crisis decision rights and escalation with pre-agreed authority
			\item Fewer approvals
			\item Lower budget
		\end{enumerate}
	\end{frame}
	
	\begin{frame}{Q191: Explanation}
		\begin{block}{Correct Answer: B}
			\textbf{Crisis authority must be predefined:}
			\begin{itemize}
				\item Crises need clear authority
				\item Predefined rights speed response
				\item More meetings slow it further
			\end{itemize}
		\end{block}
		\begin{alertblock}{Critical Trap}
			Crisis authority must be predefined.
		\end{alertblock}
	\end{frame}
	
	\begin{frame}{Q192: Executive Performance}
		\begin{block}{Question}
			An executive is rewarded despite repeated control failures. What governance signal is most damaging?
		\end{block}
		\begin{enumerate}[(A)]
			\item Incentives matter more than risk
			\item Strong oversight
			\item Higher morale
			\item Better governance
		\end{enumerate}
	\end{frame}
	
	\begin{frame}{Q192: Explanation}
		\begin{block}{Correct Answer: A}
			\textbf{Rewards reveal true priorities:}
			\begin{itemize}
				\item Rewards signal priorities
				\item Rewarding failure undermines risk culture
				\item Trust erodes
			\end{itemize}
		\end{block}
		\begin{alertblock}{Critical Trap}
			Rewards reveal true priorities.
		\end{alertblock}
	\end{frame}
	
	\begin{frame}{Q193: Board Reporting Bias}
		\begin{block}{Question}
			Reports to the board systematically minimize risk. What governance outcome is MOST likely?
		\end{block}
		\begin{enumerate}[(A)]
			\item Strong oversight
			\item Poor decisions and weakened oversight
			\item Lower risk
			\item Better culture
		\end{enumerate}
	\end{frame}
	
	\begin{frame}{Q193: Explanation}
		\begin{block}{Correct Answer: B}
			\textbf{Minimized risk leads to missed signals:}
			\begin{itemize}
				\item Boards need honest input
				\item Bias leads to missed signals
				\item Trust is fragile
			\end{itemize}
		\end{block}
		\begin{alertblock}{Critical Trap}
			Minimized risk leads to missed signals.
		\end{alertblock}
	\end{frame}
	
	\begin{frame}{Q194: Governance Refresh}
		\begin{block}{Question}
			The governance framework hasn't been reviewed in years. What should governance do FIRST?
		\end{block}
		\begin{enumerate}[(A)]
			\item Add new members
			\item Assess current relevance and align with strategy and regulation
			\item Increase meetings
			\item Reduce reporting
		\end{enumerate}
	\end{frame}
	
	\begin{frame}{Q194: Explanation}
		\begin{block}{Correct Answer: B}
			\textbf{Frameworks must evolve:}
			\begin{itemize}
				\item Frameworks can drift
				\item Alignment with strategy is essential
				\item Refresh improves effectiveness
			\end{itemize}
		\end{block}
		\begin{alertblock}{Critical Trap}
			Frameworks must evolve with the business.
		\end{alertblock}
	\end{frame}
	
	\begin{frame}{Q195: Business Case Governance}
		\begin{block}{Question}
			A large program bypasses business-case review. What governance improvement is MOST valuable?
		\end{block}
		\begin{enumerate}[(A)]
			\item More auditors
			\item Defined approval gates and criteria for large programs
			\item Reduce reporting
			\item Add more vendors
		\end{enumerate}
	\end{frame}
	
	\begin{frame}{Q195: Explanation}
		\begin{block}{Correct Answer: B}
			\textbf{Gates and criteria enforce review:}
			\begin{itemize}
				\item Gates enforce review
				\item Criteria ensure consistency
				\item More auditors do not fix design
			\end{itemize}
		\end{block}
		\begin{alertblock}{Critical Trap}
			Gates and criteria enforce review.
		\end{alertblock}
	\end{frame}
	
	\begin{frame}{Q196: Whistleblower Response}
		\begin{block}{Question}
			A whistleblower reports wrongdoing but receives no response. What is the governance impact?
		\end{block}
		\begin{enumerate}[(A)]
			\item Nothing
			\item Loss of trust and future silence
			\item Faster reporting
			\item Better culture
		\end{enumerate}
	\end{frame}
	
	\begin{frame}{Q196: Explanation}
		\begin{block}{Correct Answer: B}
			\textbf{Silent response means future silence:}
			\begin{itemize}
				\item Silence discourages reporting
				\item Trust decays
				\item Governance requires response
			\end{itemize}
		\end{block}
		\begin{alertblock}{Critical Trap}
			Silent response means future silence.
		\end{alertblock}
	\end{frame}
	
	\begin{frame}{Q197: Risk Appetite Recalibration}
		\begin{block}{Question}
			The business environment is shifting rapidly. What should governance do with appetite?
		\end{block}
		\begin{enumerate}[(A)]
			\item Ignore the shift
			\item Reassess appetite against the new environment with the board
			\item Lower appetite without review
			\item Freeze all decisions
		\end{enumerate}
	\end{frame}
	
	\begin{frame}{Q197: Explanation}
		\begin{block}{Correct Answer: B}
			\textbf{Appetite must reflect the environment:}
			\begin{itemize}
				\item Appetite must fit the environment
				\item Board re-approval ensures alignment
				\item Freezing decisions stalls the business
			\end{itemize}
		\end{block}
		\begin{alertblock}{Critical Trap}
			Appetite must reflect the environment.
		\end{alertblock}
	\end{frame}
	
	\begin{frame}{Q198: Board Independence}
		\begin{block}{Question}
			A board member has a financial interest in a vendor under decision. What governance action is required?
		\end{block}
		\begin{enumerate}[(A)]
			\item Vote on the vendor
			\item Disclose and recuse
			\item Ignore
			\item Delegate to the CEO
		\end{enumerate}
	\end{frame}
	
	\begin{frame}{Q198: Explanation}
		\begin{block}{Correct Answer: B}
			\textbf{Disclose and recuse, always:}
			\begin{itemize}
				\item Conflicts must be disclosed
				\item Recusal protects integrity
				\item Delegation does not remove the conflict
			\end{itemize}
		\end{block}
		\begin{alertblock}{Critical Trap}
			Disclose and recuse, always.
		\end{alertblock}
	\end{frame}
	
	\begin{frame}{Q199: Performance Escalation}
		\begin{block}{Question}
			Repeated missed control commitments from a business unit. What is governance's BEST action?
		\end{block}
		\begin{enumerate}[(A)]
			\item Ignore
			\item Require corrective action with accountability and reporting
			\item Increase budget
			\item Reduce monitoring
		\end{enumerate}
	\end{frame}
	
	\begin{frame}{Q199: Explanation}
		\begin{block}{Correct Answer: B}
			\textbf{Repeated misses demand intervention:}
			\begin{itemize}
				\item Missed commitments need action
				\item Accountability drives change
				\item Ignoring normalizes failure
			\end{itemize}
		\end{block}
		\begin{alertblock}{Critical Trap}
			Repeated misses demand intervention.
		\end{alertblock}
	\end{frame}
	
	\begin{frame}{Q200: Ultimate Governance}
		\begin{block}{Question}
			When board oversight is weak, what is the MOST likely consequence?
		\end{block}
		\begin{enumerate}[(A)]
			\item Strong risk management
			\item Unmanaged risk and stakeholder harm
			\item Better controls
			\item Lower exposure
		\end{enumerate}
	\end{frame}
	
	\begin{frame}{Q200: Explanation}
		\begin{block}{Correct Answer: B}
			\textbf{Weak oversight exposes stakeholders:}
			\begin{itemize}
				\item Oversight is a foundation of governance
				\item Weak boards invite drift
				\item Stakeholders bear the consequences
			\end{itemize}
		\end{block}
		\begin{alertblock}{Critical Trap}
			Weak oversight exposes stakeholders.
		\end{alertblock}
	\end{frame}
	
	\begin{frame}{Exam Mindset Summary}
		\begin{block}{Remember}
			\begin{itemize}
				\item \textbf{Governance} directs; \textbf{management} executes
				\item \textbf{Accountability} cannot be delegated or outsourced
				\item \textbf{Appetite} is board-owned; \textbf{tolerance} is permitted variation
				\item \textbf{Three lines:} operate (1st), oversee (2nd), audit (3rd)
				\item \textbf{Policy} states intent; \textbf{standards} specify requirements
				\item \textbf{Culture} follows leadership behavior, not slogans
				\item \textbf{Reporting} needs honesty, clarity, and escalation paths
				\item \textbf{Third parties} need tiering, exit plans, and oversight
			\end{itemize}
		\end{block}
	\end{frame}
	
	\begin{frame}{Thank You}
		\begin{center}
			{\Huge \textbf{Good Luck on Your CRISC Exam!}}
		\end{center}
	\end{frame}
	
\end{document}